\documentclass[12pt,a4paper]{article}
\usepackage{amsmath,amssymb,amsthm}
\usepackage{hyperref}
\usepackage[margin=1in]{geometry}

\newtheorem{theorem}{Theorem}[section]
\newtheorem{proposition}[theorem]{Proposition}
\newtheorem{lemma}[theorem]{Lemma}
\newtheorem{corollary}[theorem]{Corollary}
\newtheorem{definition}[theorem]{Definition}
\newtheorem{remark}[theorem]{Remark}

\newcommand{\HH}{\mathrm{HH}}
\newcommand{\Hom}{\mathrm{Hom}}
\newcommand{\kk}{\mathbf{k}}
\newcommand{\CC}{\mathbb{C}}
\newcommand{\ZZ}{\mathbb{Z}}
\newcommand{\im}{\mathrm{im}}

\title{Supplementary Material: Proof of the Decomposition Theorem}
\author{L.\ Horesh \and Claude}
\date{2026}

\begin{document}
\maketitle

\section{Proof of Theorem~4.1 (Main Decomposition)}
\label{sec:proof-decomposition}

We prove:
\[
\HH^2(A_M, A_M) \cong H^2(G_M, \CC) \oplus V_{\mathrm{phonon}}(M).
\]

\subsection{Step 1: The pullback map}

Let $G = G_M$ be the crystallographic point group of material $M$,
and let $A = A_M$ be the bond algebra. The group $G$ acts on $A$ by
algebra automorphisms: for $g \in G$ and an arrow $\alpha: i \to j$
in the bond quiver, $g \cdot \alpha = \alpha_{g(i), g(j)}$ (the
arrow between the permuted atoms).

\begin{definition}
The \emph{pullback map} $\iota^*: C^n(G, \CC) \to C^n(A, A)$ is
defined as follows. For a group $n$-cochain $f: G^n \to \CC$, define
the Hochschild $n$-cochain $\iota^*(f): A^{\otimes n} \to A$ by:
\[
\iota^*(f)(a_1 \otimes \cdots \otimes a_n) =
f(g_{a_1}, \ldots, g_{a_n}) \cdot (a_1 \cdots a_n)
\]
where $g_{a_i} \in G$ is the group element mapping the source vertex
of the path $a_i$ to its target vertex (if $a_i$ is an arrow), and
extended by linearity and the Leibniz rule for longer paths.
\end{definition}

\begin{lemma}[Compatibility with differentials]
\label{lem:pullback-cochain-map}
$\iota^*$ is a cochain map: $d^n_{\HH} \circ \iota^* = \iota^* \circ d^n_G$.
\end{lemma}

\begin{proof}
Direct computation. For a group 2-cocycle $f: G \times G \to \CC$
satisfying $f(g,h) + f(gh,k) = f(g,hk) + f(h,k)$, we verify that
$\iota^*(f)$ satisfies the Hochschild 2-cocycle condition:
\begin{align*}
a_1 \cdot \iota^*(f)(a_2, a_3)
&+ \iota^*(f)(a_1 a_2, a_3) - \iota^*(f)(a_1, a_2 a_3) - \iota^*(f)(a_1, a_2) \cdot a_3 \\
&= f(g_1, g_2) \cdot a_1 a_2 a_3 + f(g_1 g_2, g_3) \cdot a_1 a_2 a_3 \\
&\quad - f(g_1, g_2 g_3) \cdot a_1 a_2 a_3 - f(g_1, g_2) \cdot a_1 a_2 a_3 \\
&= [f(g_1 g_2, g_3) - f(g_1, g_2 g_3)] \cdot a_1 a_2 a_3 \\
&= [f(g_1, g_2) - f(g_2, g_3)] \cdot a_1 a_2 a_3 = 0
\end{align*}
by the group cocycle condition.
\end{proof}

\subsection{Step 2: Injectivity of the pullback}

\begin{proposition}
\label{prop:pullback-injective}
The induced map $\iota^*: H^2(G, \CC) \to \HH^2(A, A)$ is injective.
\end{proposition}

\begin{proof}
Suppose $\iota^*(f) = d^1_{\HH}(\phi)$ for some $\phi \in C^1(A, A)$.
We construct a group 1-cochain $\psi: G \to \CC$ such that
$f = d^1_G(\psi)$, contradicting $[f] \neq 0$ in $H^2(G, \CC)$.

Define $\psi(g) = \frac{1}{|G|} \sum_{h \in G} \phi(e_{h(v_0)})$
where $v_0$ is a fixed base vertex and $e_v$ denotes the idempotent
at vertex $v$. This is well-defined because $\CC$ has characteristic
zero and $|G|$ is invertible.

The averaging argument shows:
\[
d^1_G(\psi)(g, h) = \psi(g) + \psi(h) - \psi(gh)
= \frac{1}{|G|} \sum_{k \in G}
[\phi(e_{k(v_0)}) + \phi(e_{hk(v_0)}) - \phi(e_{ghk(v_0)})]
= f(g, h)
\]
where the last equality uses $d^1_{\HH}(\phi) = \iota^*(f)$ evaluated
on appropriate algebra elements, and averaging over orbits.
\end{proof}

\subsection{Step 3: The orthogonal complement}

\begin{definition}
The \emph{trace inner product} on $C^n(A, A)$ is:
\[
\langle f, g \rangle = \sum_{\text{basis paths } p}
\overline{f(p)} \cdot g(p)
\]
where the sum runs over a basis of $A^{\otimes n}$.
\end{definition}

\begin{definition}
The \emph{phonon-deformation space} is:
\[
V_{\mathrm{phonon}}(M) := \{ [\alpha] \in \HH^2(A_M) :
\langle \alpha, \iota^*(\beta) \rangle = 0
\text{ for all } [\beta] \in H^2(G, \CC) \}
\]
\end{definition}

\begin{theorem}[Main Decomposition]
\[
\HH^2(A_M, A_M) = \im(\iota^*) \oplus V_{\mathrm{phonon}}(M)
\cong H^2(G_M, \CC) \oplus V_{\mathrm{phonon}}(M).
\]
\end{theorem}

\begin{proof}
By Proposition~\ref{prop:pullback-injective}, $\im(\iota^*)$ is a
subspace isomorphic to $H^2(G, \CC)$. Since $\HH^2(A_M)$ is
finite-dimensional (the bond quiver is finite), the trace inner
product is non-degenerate, and $V_{\mathrm{phonon}}$ is indeed the
orthogonal complement. The direct sum decomposition follows.
\end{proof}

\section{Dimension formula for $V_{\mathrm{phonon}}$}

\begin{corollary}
$\dim V_{\mathrm{phonon}}(M) = \dim \HH^2(A_M) - \dim H^2(G_M, \CC)$.
\end{corollary}

For crystallographic point groups, $\dim H^2(G, \CC)$ is given by the
Schur multiplier table:
\begin{center}
\begin{tabular}{lcc}
\hline
Point group & $|G|$ & $\dim H^2(G, \CC)$ \\
\hline
$C_1$ & 1 & 0 \\
$C_i, C_s, C_2$ & 2 & 0 \\
$C_{2v}, C_{2h}, S_4$ & 4 & 0 \\
$D_2 \cong V_4$ & 4 & 1 \\
$D_{2h}$ & 8 & 3 \\
$D_{3h}$ & 12 & 0 \\
$D_{4h}$ & 16 & 2 \\
$T$ & 12 & 1 \\
$T_d$ & 24 & 1 \\
$O$ & 24 & 1 \\
$O_h$ & 48 & 1 \\
\hline
\end{tabular}
\end{center}

Thus for materials with $O_h$ symmetry (e.g., cubic perovskites),
$\dim V_{\mathrm{phonon}} = \dim \HH^2 - 1$, and most of HH^2
content is phonon-deformation.

\section{Physical interpretation of $V_{\mathrm{phonon}}$}

Each linearly independent element $[\alpha] \in V_{\mathrm{phonon}}$
represents an infinitesimal deformation of the bond algebra that is:
\begin{enumerate}
\item Not induced by the point-group symmetry structure;
\item Compatible with the triangle-commutativity relations (it
  deforms the algebra \emph{within} the class of bond algebras);
\item A \emph{soft mode}: it costs zero energy at the infinitesimal
  level (the deformation is unobstructed in $\HH^2$).
\end{enumerate}

The connection to phonon coupling is:
\begin{itemize}
\item Each soft-mode direction in $V_{\mathrm{phonon}}$ corresponds
  to a specific pattern of bond-length changes.
\item These bond-length changes couple to electronic states via the
  deformation potential.
\item Materials with more such modes (larger $\dim V_{\mathrm{phonon}}$)
  have more channels for electron-phonon coupling.
\item This explains the observed correlation with $\lambda$.
\end{itemize}

\end{document}
